% !TEX program = xelatex
\documentclass[11pt,a4paper]{article}

\usepackage[UTF8]{ctex}
\usepackage[margin=2.2cm]{geometry}
\usepackage{amsmath, amssymb, amsthm}
\usepackage{listings}
\usepackage{xcolor}
\usepackage{tikz}
\usepackage{booktabs}
\usepackage{multirow}
\usepackage{hyperref}
\usetikzlibrary{positioning, arrows.meta, shapes.geometric, fit, calc, decorations.pathreplacing}

\hypersetup{
  colorlinks=true,
  linkcolor=blue!60!black,
  urlcolor=teal!70!black,
}

\lstdefinelanguage{Rust}{
  morekeywords={fn,let,mut,pub,struct,impl,use,mod,self,Self,Option,Some,None,return,if,else,for,while,loop,match,break,continue,as,where,trait,type,enum,move,ref,const,static,unsafe,extern,crate,super,dyn,async,await,box,in},
  sensitive=true,
  morecomment=[l]{//},
  morecomment=[s]{/*}{*/},
  morestring=[b]",
}

\lstset{
  language=Rust,
  basicstyle=\ttfamily\footnotesize,
  keywordstyle=\color{blue!60!black}\bfseries,
  commentstyle=\color{green!40!black}\itshape,
  stringstyle=\color{red!60!black},
  numbers=left,
  numberstyle=\tiny\color{gray},
  numbersep=8pt,
  frame=single,
  framerule=0.4pt,
  rulecolor=\color{gray!40},
  breaklines=true,
  breakatwhitespace=true,
  showstringspaces=false,
  tabsize=4,
  columns=fullflexible,
  keepspaces=true,
}

\title{\textbf{halfedge\,: BVH 加速结构模块设计文档} \\ \large \texttt{src/bvh.rs}}
\author{halfedge 库}
\date{2026-07-05}

\begin{document}
\maketitle
\tableofcontents

\section{模块定位}

\texttt{bvh.rs} 在 \texttt{geometry.rs} 之上叠加\textbf{空间加速结构}能力，
将射线求交与最近点查询从 $O(F)$ 暴力扫描降到平均 $O(\log F)$。
对大网格（$> 10^3$ 面）有数量级加速，对小网格（$< 64$ 面）BVH 自身常数
开销可能让它略慢于暴力。

\textbf{核心 API}：

\begin{center}
\renewcommand{\arraystretch}{1.18}
\begin{tabular}{@{}lll@{}}
  \toprule
  \textbf{类别} & \textbf{API} & \textbf{语义} \\
  \midrule
  \multirow{2}{*}{构建}
    & \texttt{Bvh::new()}                  & 创建空 BVH \\
    & \texttt{Bvh::build(\&mesh)}          & 从网格构建（仅三角面），返回 \texttt{Bvh} \\
  \midrule
  \multirow{2}{*}{查询}
    & \texttt{bvh.ray\_intersection(o, d, \&mesh)} & 射线求最近交点，返回 \texttt{Option<RayHit>} \\
    & \texttt{bvh.nearest\_point(p, \&mesh)}      & 网格表面最近点，返回 \texttt{(FaceId, [f64;3], f64)} \\
  \midrule
  \multirow{3}{*}{统计}
    & \texttt{bvh.node\_count() / leaf\_count() / depth()} & 节点数 / 叶子数 / 树深 \\
    & \texttt{bvh.is\_empty()}                                & 是否无三角面 \\
    & \texttt{LEAF\_MAX\_FACES}（常量 = 8）                   & 叶子节点最大面数 \\
  \bottomrule
\end{tabular}
\end{center}

\section{数据结构}

BVH 是一棵二叉树，节点统一表示为 \texttt{BvhNode}：

\begin{lstlisting}
struct BvhNode {
    aabb: AABB,                 // 节点合并包围盒
    left: Option<usize>,        // 内部节点：左子索引
    right: Option<usize>,       // 内部节点：右子索引
    faces: Vec<FaceId>,         // 叶子节点：面 ID 列表
}
\end{lstlisting}

\begin{itemize}
  \item \textbf{内部节点}：\texttt{left} 与 \texttt{right} 都为 \texttt{Some}，
        \texttt{faces} 为空；
  \item \textbf{叶子节点}：\texttt{left} 与 \texttt{right} 都为 \texttt{None}，
        \texttt{faces} 非空；
  \item \texttt{is\_leaf()} 通过 \texttt{left.is\_none() \&\& right.is\_none()} 判定。
\end{itemize}

\texttt{Bvh} 持有 \texttt{nodes: Vec<BvhNode>} 与 \texttt{root: Option<usize>}，
全部节点存储在连续 \texttt{Vec} 中（cache-friendly）。

\section{构建算法}

\subsection{整体流程}

\begin{center}
\resizebox{\textwidth}{!}{%
\begin{tikzpicture}[
  node distance=0.5cm,
  proc/.style={draw, rounded corners=2pt, minimum width=11cm, minimum height=0.7cm, align=center, font=\small},
  dec/.style={diamond, draw, aspect=2.4, fill=orange!12, inner sep=1pt, font=\small, align=center, minimum width=3.5cm},
  op/.style={-Stealth, thick},
  startend/.style={draw, rounded corners=10pt, minimum width=2.6cm, minimum height=0.7cm, font=\small, fill=gray!12}
]
  \node[startend] (s) {开始};
  \node[proc, below=of s, fill=blue!8] (p1) {1. 收集每面 \texttt{(FaceId, AABB, center)}（跳过非三角 / 退化）};
  \node[dec, below=of p1] (d1) {prims 为空？};
  \node[proc, right=2.0cm of d1, fill=red!10] (empty) {返回空 \texttt{Bvh}};
  \node[proc, below=of d1, fill=blue!8] (p2) {2. 调用 \texttt{build\_recursive(prims, nodes)}};
  \node[startend, below=of p2] (e) {返回 \texttt{Bvh}};
  \draw[op] (s) -- (p1);
  \draw[op] (p1) -- (d1);
  \draw[op] (d1) -- node[above, font=\scriptsize] {是} (empty);
  \draw[op] (d1) -- node[right, font=\scriptsize] {否} (p2);
  \draw[op] (p2) -- (e);
\end{tikzpicture}
}
\end{center}

\subsection{递归分裂}

\texttt{build\_recursive} 对当前 \texttt{prims} 列表：

\begin{enumerate}
  \item 计算 \textbf{节点 AABB}：所有 primitive AABB 的合并；
  \item \textbf{终止条件}：若 $|\text{prims}| \le \texttt{LEAF\_MAX\_FACES} = 8$，
        创建叶子节点（保存面 ID 列表）并返回；
  \item \textbf{选择最长轴}：计算 AABB 对角线 $\vec{d}$，取
        $\text{axis} = \arg\max_i d_i$；
  \item \textbf{按面中心排序}：\texttt{prims.sort\_by} 按 $\text{center[axis]}$ 升序；
  \item \textbf{中位数划分}：\texttt{mid = prims.len() / 2}，
        \texttt{right\_prims = prims.split\_off(mid)}；
  \item \textbf{占位父节点}：先 \texttt{push} 父节点（\texttt{left/right = None}），
        再递归构建左右子树，最后回填索引。
\end{enumerate}

\begin{center}
\begin{tikzpicture}[
  node distance=0.55cm,
  proc/.style={draw, rounded corners=2pt, minimum width=10cm, minimum height=0.7cm, align=center, font=\small},
  op/.style={-Stealth, thick},
  startend/.style={draw, rounded corners=10pt, minimum width=2.6cm, minimum height=0.7cm, font=\small, fill=gray!12}
]
  \node[startend] (s) {进入 build\_recursive};
  \node[proc, below=of s, fill=blue!8] (p1) {计算节点 AABB（合并所有 primitive）};
  \node[proc, below=of p1, fill=yellow!15] (p2) {$|\text{prims}| \le 8$？创建叶子返回};
  \node[proc, below=of p2, fill=blue!8] (p3) {选最长轴 axis $= \arg\max_i d_i$};
  \node[proc, below=of p3, fill=blue!8] (p4) {按 center[axis] 排序 prims};
  \node[proc, below=of p4, fill=blue!8] (p5) {mid $= |\text{prims}| / 2$；right $\gets$ prims.split\_off(mid)};
  \node[proc, below=of p5, fill=green!12] (p6) {占位父节点；递归 left $=$ build\_recursive(prims)};
  \node[proc, below=of p6, fill=green!12] (p7) {递归 right $=$ build\_recursive(right\_prims)};
  \node[proc, below=of p7, fill=green!18] (p8) {回填 nodes[idx].left / right};
  \node[startend, below=of p8] (e) {返回 idx};
  \draw[op] (s) -- (p1);
  \draw[op] (p1) -- (p2);
  \draw[op] (p2) -- (p3);
  \draw[op] (p3) -- (p4);
  \draw[op] (p4) -- (p5);
  \draw[op] (p5) -- (p6);
  \draw[op] (p6) -- (p7);
  \draw[op] (p7) -- (p8);
  \draw[op] (p8) -- (e);
\end{tikzpicture}
\end{center}

\subsection{退化守卫}

构建时跳过：
\begin{itemize}
  \item 非三角面（边界环长度 $\ne 3$）；
  \item 顶点缺失的面；
  \item 退化三角形（AABB 对角线 $< 10^{-14}$，含重复顶点或共线顶点）。
\end{itemize}

\subsection{复杂度}

\begin{itemize}
  \item 时间：$O(F \log F)$（每层 $O(F)$ 排序，共 $\log F$ 层）；
  \item 空间：$O(F)$（节点数 $\le 2F - 1$，满二叉树上限）；
  \item 树深：平均 $O(\log F)$，最坏 $O(F)$（极端退化输入）。
\end{itemize}

\section{射线-AABB 求交：Slab 法}

\subsection{公式}

AABB 由 \texttt{min} / \texttt{max} 三对值定义。对每个轴 $i \in \{x, y, z\}$，
slab 范围是 $[\text{min}_i, \text{max}_i]$。射线 $\vec{R}(t) = \vec{O} + t\vec{D}$
在轴 $i$ 上的进入/退出参数：
\[
  t_i^{\min} = \frac{\text{min}_i - O_i}{D_i}, \quad
  t_i^{\max} = \frac{\text{max}_i - O_i}{D_i},
\]
若 $t_i^{\min} > t_i^{\max}$ 互换。

\textbf{相交条件}：
\[
  t_{\text{enter}} = \max_i \min(t_i^{\min}, t_i^{\max}), \quad
  t_{\text{exit}}  = \min_i \max(t_i^{\min}, t_i^{\max}),
\]
\[
  t_{\text{enter}} \le t_{\text{exit}} \quad \wedge \quad t_{\text{exit}} \ge 0.
\]

\subsection{平行 slab 处理}

若 $|D_i| < 10^{-14}$（射线与 slab 平行），则需 $O_i \in [\text{min}_i, \text{max}_i]$，
否则不相交。

\subsection{返回值}

返回 \texttt{Option<(t\_enter, t\_exit)>}：
\begin{itemize}
  \item \texttt{t\_enter}：射线进入 AABB 的参数（可为负，原点在 AABB 内时）；
  \item \texttt{t\_exit}：射线离开 AABB 的参数。
\end{itemize}

\section{射线-网格求交（BVH 加速）}

\subsection{递归算法}

\begin{center}
\resizebox{\textwidth}{!}{%
\begin{tikzpicture}[
  node distance=0.5cm,
  proc/.style={draw, rounded corners=2pt, minimum width=11cm, minimum height=0.7cm, align=center, font=\small},
  dec/.style={diamond, draw, aspect=2.4, fill=orange!12, inner sep=1pt, font=\small, align=center, minimum width=3.5cm},
  op/.style={-Stealth, thick},
  startend/.style={draw, rounded corners=10pt, minimum width=2.6cm, minimum height=0.7cm, font=\small, fill=gray!12}
]
  \node[startend] (s) {进入 ray\_recursive(node)};
  \node[dec, below=of s] (d1) {ray\_aabb 命中？};
  \node[proc, right=2.0cm of d1, fill=red!10] (miss) {剪枝返回};
  \node[dec, below=of d1] (d2) {$t_{\text{enter}} > \text{best\_t}$？};
  \node[proc, right=2.0cm of d2, fill=red!10] (prune) {剪枝返回};
  \node[dec, below=of d2] (d3) {是叶子？};
  \node[proc, right=2.0cm of d3, fill=green!12] (leaf) {对每个面调用 M\"oller-Trumbore；更新 best};
  \node[proc, below=of d3, fill=blue!8] (inner) {选最近 AABB 优先；递归 first / second};
  \node[startend, below=of inner] (e) {返回};
  \draw[op] (s) -- (d1);
  \draw[op] (d1) -- node[above, font=\scriptsize] {否} (miss);
  \draw[op] (d1) -- node[right, font=\scriptsize] {是} (d2);
  \draw[op] (d2) -- node[above, font=\scriptsize] {是} (prune);
  \draw[op] (d2) -- node[right, font=\scriptsize] {否} (d3);
  \draw[op] (d3) -- node[above, font=\scriptsize] {是} (leaf);
  \draw[op] (d3) -- node[right, font=\scriptsize] {否} (inner);
  \draw[op] (inner) -- (e);
\end{tikzpicture}
}
\end{center}

\subsection{剪枝条件}

\begin{enumerate}
  \item \textbf{AABB 剪枝}：射线与节点 AABB 不相交 → 整个子树跳过；
  \item \textbf{距离剪枝}：当前已知最近交点 \texttt{best\_t} 小于 AABB 进入参数
        \texttt{t\_enter} → 子树不可能产生更近交点；
  \item \textbf{叶子求交}：对叶子节点内的每个面调用
        \texttt{ray\_triangle\_intersection}（M\"oller-Trumbore），
        更新 \texttt{best} 为 $t$ 最小者。
\end{enumerate}

\subsection{子节点访问顺序}

内部节点的两个子节点都先计算 \texttt{ray\_aabb} 的 \texttt{t\_enter}，
\textbf{优先访问 t 较小者}。这样后续子树访问时 \texttt{best\_t} 已被更新为更紧的上界，
剪枝效率更高。

\subsection{复杂度}

平均 $O(\log F)$；最坏 $O(F)$（射线擦过大量 AABB）。

\section{最近点查询（BVH 加速）}

\subsection{点-AABB 最近距离平方}

对查询点 $\vec{P}$ 与 AABB：
\[
  d^2(\vec{P}, \text{AABB}) = \sum_{i} \max(0, \text{min}_i - P_i)^2 + \max(0, P_i - \text{max}_i)^2.
\]
点在 AABB 内时距离为 0。

\subsection{递归算法}

\begin{center}
\resizebox{\textwidth}{!}{%
\begin{tikzpicture}[
  node distance=0.5cm,
  proc/.style={draw, rounded corners=2pt, minimum width=11cm, minimum height=0.7cm, align=center, font=\small},
  dec/.style={diamond, draw, aspect=2.4, fill=orange!12, inner sep=1pt, font=\small, align=center, minimum width=3.5cm},
  op/.style={-Stealth, thick},
  startend/.style={draw, rounded corners=10pt, minimum width=2.6cm, minimum height=0.7cm, font=\small, fill=gray!12}
]
  \node[startend] (s) {进入 nearest\_recursive(node)};
  \node[proc, below=of s, fill=blue!8] (p1) {计算 $d^2_{\text{aabb}} = $ point\_aabb\_distance\_sq};
  \node[dec, below=of p1] (d1) {$d^2_{\text{aabb}} > \text{best\_d}^2$？};
  \node[proc, right=2.0cm of d1, fill=red!10] (prune) {剪枝返回};
  \node[dec, below=of d1] (d2) {是叶子？};
  \node[proc, right=2.0cm of d2, fill=green!12] (leaf) {对每个面调用 closest\_point\_on\_triangle；更新 best};
  \node[proc, below=of d2, fill=blue!8] (inner) {按 AABB 距离升序访问子节点};
  \node[startend, below=of inner] (e) {返回};
  \draw[op] (s) -- (p1);
  \draw[op] (p1) -- (d1);
  \draw[op] (d1) -- node[above, font=\scriptsize] {是} (prune);
  \draw[op] (d1) -- node[right, font=\scriptsize] {否} (d2);
  \draw[op] (d2) -- node[above, font=\scriptsize] {是} (leaf);
  \draw[op] (d2) -- node[right, font=\scriptsize] {否} (inner);
  \draw[op] (inner) -- (e);
\end{tikzpicture}
}
\end{center}

\subsection{剪枝条件}

\textbf{AABB 距离剪枝}：若节点 AABB 的最近距离平方
$d^2_{\text{aabb}}$ 大于当前已知最近距离平方 $\text{best\_d}^2$，
则该子树不可能产生更近的交点 → 剪枝。

\textbf{子节点访问顺序}：内部节点的两个子节点都先计算 \texttt{point\_aabb\_distance\_sq}，
\textbf{按距离升序访问}。先访问近的子树，使 \texttt{best\_d} 更紧，剪枝更高效。

\subsection{复杂度}

平均 $O(\log F)$；最坏 $O(F)$（查询点周围面密度极高）。

\section{与暴力算法的关系}

\begin{center}
\renewcommand{\arraystretch}{1.18}
\begin{tabular}{@{}llll@{}}
  \toprule
  \textbf{操作} & \textbf{暴力（geometry.rs）} & \textbf{BVH（bvh.rs）} & \textbf{加速比（大网格）} \\
  \midrule
  射线求交   & \texttt{ray\_mesh\_intersection} $O(F)$ & \texttt{bvh.ray\_intersection} $O(\log F)$ & $\sim F / \log F$ \\
  最近点     & 遍历面 + \texttt{point\_triangle\_distance} $O(F)$ & \texttt{bvh.nearest\_point} $O(\log F)$ & $\sim F / \log F$ \\
  \bottomrule
\end{tabular}
\end{center}

\textbf{何时用 BVH}：
\begin{itemize}
  \item 网格 $> 10^3$ 面，且需多次射线 / 最近点查询；
  \item BVH 构建开销 $O(F \log F)$，单次查询摊销下值得。
\end{itemize}

\textbf{何时用暴力}：
\begin{itemize}
  \item 网格 $< 64$ 面，BVH 常数开销不划算；
  \item 仅需一次或几次查询，构建开销无法摊销。
\end{itemize}

\section{退化守卫汇总}

\begin{center}
\renewcommand{\arraystretch}{1.18}
\begin{tabular}{@{}lll@{}}
  \toprule
  \textbf{函数} & \textbf{退化场景} & \textbf{处理} \\
  \midrule
  \texttt{Bvh::build}            & 空网格 / 无三角面 & 返回空 \texttt{Bvh}（\texttt{root = None}） \\
  \texttt{Bvh::build}            & 退化三角形（共线 / 重复顶点） & 跳过该面 \\
  \texttt{ray\_intersection}     & 空 BVH & 返回 \texttt{None} \\
  \texttt{nearest\_point}        & 空 BVH & 返回 \texttt{None} \\
  \texttt{ray\_aabb}             & 射线与 slab 平行 & 检查原点是否在 slab 内 \\
  \texttt{ray\_aabb}             & AABB 在射线反方向 & \texttt{t\_exit < 0} 返回 \texttt{None} \\
  \texttt{point\_aabb\_distance\_sq} & 点在 AABB 内 & 返回 0 \\
  \bottomrule
\end{tabular}
\end{center}

\section{测试覆盖}

\begin{center}
\renewcommand{\arraystretch}{1.18}
\begin{tabular}{@{}ll@{}}
  \toprule
  \textbf{测试名} & \textbf{验证点} \\
  \midrule
  \multicolumn{2}{l}{\textit{构建}} \\
  \texttt{build\_empty\_mesh}                         & 空网格 BVH 为空 \\
  \texttt{build\_two\_triangles\_basic}               & 2 面 $\le 8$ → 单叶子节点 \\
  \texttt{large\_mesh\_bvh\_depth\_is\_logarithmic}   & icosphere(3) 80 面 BVH 深度 $\le 12$ \\
  \midrule
  \multicolumn{2}{l}{\textit{射线求交}} \\
  \texttt{ray\_hit\_first\_triangle}                  & 射线命中第一个三角形 \\
  \texttt{ray\_hit\_second\_triangle\_when\_first\_missed} & 射线只命中第二个三角形 \\
  \texttt{ray\_miss\_returns\_none}                   & 远离网格的射线返回 \texttt{None} \\
  \texttt{ray\_returns\_nearest\_hit}                 & icosphere 球面命中半径 $\approx 1$ \\
  \texttt{icosphere\_bvh\_consistent\_with\_brute\_force\_ray} & BVH 与暴力算法 $t$ 一致 \\
  \midrule
  \multicolumn{2}{l}{\textit{最近点}} \\
  \texttt{nearest\_point\_on\_triangle\_returns\_vertex} & 远点最近点为顶点 \\
  \texttt{nearest\_point\_inside\_triangle}           & 三角形上方点投影到面 \\
  \texttt{nearest\_point\_far\_away\_picks\_closer\_triangle} & 选择更近的三角形 \\
  \texttt{icosphere\_bvh\_consistent\_with\_brute\_force\_nearest} & BVH 与暴力最近距离一致 \\
  \midrule
  \multicolumn{2}{l}{\textit{辅助函数}} \\
  \texttt{ray\_aabb\_basic\_hit}                      & 射线-AABB 命中 / 未命中 \\
  \texttt{ray\_aabb\_parallel\_slab}                  & 平行 slab 边界条件 \\
  \texttt{point\_aabb\_distance\_basic}               & 点-AABB 距离（内部 / 角落 / 面外） \\
  \bottomrule
\end{tabular}
\end{center}

\texttt{src/bvh.rs} 共 \textbf{15 个}单元测试，全库共 \textbf{371 个}。

\end{document}
